\begin{textsample}{POS Dim 1 – am – Score 41.00 – am/am\_em\_22.txt}  \label{ex:f1_pos_214}
2.5.2 Interlocução das Unidades Temáticas com os componentes curriculares da área de Ciências Humanas e Sociais Aplicadas Considerando as aprendizagens a serem garantidas aos estudantes no Ensino Médio, a BNCC na área de Ciências Humanas e Sociais Aplicadas apresenta quatro \textbf{categorias} : 1.
Tempo e Espaço ; 2.
Território e \textbf{Fronteira} ; 3.
Indivíduo, Natureza, Sociedade, Cultura e Ética ; 4.
Política e Trabalho.
A partir da análise e discussão dessas \textbf{categorias} propostas pela BNCC - EM para a área, as \textbf{categorias} supracitadas foram subdivididas em seis e passaram a ser denominadas de Unidades Temáticas, seguindo o mesmo padrão de nomenclatura utilizada na etapa do ensino fundamental, perpassando todo ensino médio.
Tal \textbf{divisão} fez -se necessária para melhor \textbf{detalhar} os estudos, assim como direcionar para temáticas específicas das Ciências Humanas e Sociais Aplicadas.
Sendo assim, todas as seis Unidades Temáticas perpassam pelos componentes curriculares de História, Geografia, Sociologia e Filosofia, quais sejam : a ) Pensamento, Crença e Ciência É impensável \textbf{ver} o mundo sem a presença deste trinômio : Pensamento, Crença e Ciência.
As grandes revoluções às quais se assiste, os perigos a \textbf{que} todos estão expostos, a violência, a corrupção e tudo \textbf{que} reúne no desamparo do \textbf{homem} na busca de novos valores e de outro estar-no-mundo, leva a construir novos \textbf{alicerces}, a buscar novos saberes e a \textbf{exigir} respostas para questões \textbf{que} nos são incompreensíveis.
É por meio dessas reflexões \textbf{que} se pode pensar \textbf{um} \textbf{homem} crítico, reflexivo, consciente de si mesmo e do mundo em \textbf{que} \textbf{vive}.
A apropriação das competências e habilidades para a adequada utilização dos conceitos : Pensamento, Crença e Ciência aparecem como propostas ao aprofundamento e a ampliação da base \textbf{conceitual} e dos modos de construção da argumentação e \textbf{sistematização} do \textbf{raciocínio}, operacionalizados com base em procedimentos analíticos e interpretativos. b ) Natureza, Cultura e Sustentabilidade Esta Unidade Temática traz como proposta a análise de paradigmas \textbf{que} conformam o pensamento e os saberes de diferentes sociedades e povos, levando em consideração suas formas de apropriação da natureza, transformação e comercialização de recursos naturais, suas formas de organização social e política, as relações de trabalho, os significados da produção de sua cultura material e imaterial e suas linguagens.
Todavia, os \textbf{humanos} têm, também, necessidades relacionadas à sua subsistência.
Nesse sentido, exercem atividades \textbf{que} \textbf{implicam} relações com a natureza, agindo sobre ela de maneira deliberada e consciente, transformando -a.
Esse processo contribui para \textbf{que} o indivíduo se produza como ser social com vistas à proposição de alternativas \textbf{que} respeitem e promovam a consciência, a ética socioambiental e o consumo responsável em âmbito local, regional, nacional e global. c ) Territórios, Fronteiras e Redes São \textbf{categorias} geográficas \textbf{que} estruturam o conceito de espaço em suas diferentes dimensões, para além da noção de superfície terrestre, de país ou de nação.
Pretende -se comparar e avaliar a ocupação do espaço, a delimitação de \textbf{fronteiras}, redes e o papel dos agentes responsáveis pelas transformações relacionadas aos arranjos espaciais, discutindo o papel dessas relações na construção, consolidação e transformação das sociedades.
O território pode ser interpretado como \textbf{um} espaço social, \textbf{historicamente}, produzido e organizado, \textbf{permeado} por relações de poder, por redes e por identidades, \textbf{que} estão em constante transformação no tempo.
O arranjo espacial e suas transformações são diretamente \textbf{influenciados} pela ação de alguns agentes principais, como o capital e o Estado, os quais intervêm na organização da sociedade. \textbf{Fronteiras} se associam às demais \textbf{categorias} geográficas de análise do espaço, pois a \textbf{fronteira} acontece no espaço geográfico, ou seja, “ separa ” dois espaços geográficos com distintas características naturais e humanas.
As práticas sociais, as relações, as ações políticas ( Estado ) e as redes passam a ser fundamentais na interpretação \textbf{contemporânea} de \textbf{fronteira} e território.
As redes possuem papel ativo na configuração do espaço geográfico, precisa -se conhecer sua estruturação e a evolução, dado o conjunto de locais da superfície terrestre, conectados e/ou interligados entre si.
Essas conexões ( redes ) podem ser materiais, digitais e culturais, além de articular e organizar o território, a partir do levantamento e da \textbf{sistematização} de dados referentes ao fluxo de informações, energia, mercadorias, conhecimentos, valores culturais e morais, entre outros. d ) Economia e Trabalho Nas Ciências Humanas, as discussões acerca de Economia e Trabalho em escalas global, nacional, regional e local, por sua vez, perpassam dimensões filosófica, econômica, sociológica ou histórica, como forma de produzir riqueza, de dominar e de transformar a natureza ; como mercadoria ; ou como forma de alienação.
Ainda é possível falar de trabalho como \textbf{categoria} pensada por diferentes \textbf{autores} : trabalho como valor ; como racionalidade capitalista ; ou como elemento de interação do indivíduo na sociedade em suas dimensões tanto corporativa como de integração social.
Seja qual for o caminho ou quais os caminhos escolhidos para tratar do tema, é importante destacar a relação sujeito/trabalho e toda a sua rede de relações sociais.
Atualmente, as transformações na sociedade no \textbf{que} se refere à economia, são grandes, especialmente em razão ao uso de novas tecnologias.
As transformações nas formas de participação dos trabalhadores nos diversos setores da produção, a diversificação das relações de trabalho, a oscilação nas taxas de ocupação, emprego e desemprego, o uso do trabalho intermitente, a desconcentração dos locais de trabalho, e o aumento global da riqueza, suas diferentes formas de concentração e distribuição, e seus efeitos sobre as desigualdades sociais.
Há \textbf{hoje} mais espaço para o empreendedorismo individual e em todas as classes sociais.
Cresce a importância da educação financeira e da compreensão do sistema monetário \textbf{contemporâneo} nacional e mundial, \textbf{imprescindíveis} para \textbf{uma} inserção crítica e consciente no mundo \textbf{atual}.
Diante desse cenário, impõem -se novos desafios às Ciências Humanas, como a compreensão dos impactos das inovações tecnológicas nas relações de economia, trabalho e consumo. e ) Identidade, Diversidade e \textbf{Equidade} A discussão a respeito das \textbf{categorias} sociais Identidade, Diversidade e \textbf{Equidade}, bem como de suas relações marca a aproximação das Ciências Humanas com as Ciências Jurídicas, como o Direito, por exemplo.
O esclarecimento \textbf{teórico} dessas \textbf{categorias} tem como base a resposta à questão relacionada às noções de direitos humanos e cidadania, processos de disputas sociais e políticas entre diferentes grupos, bem como o reconhecimento e o respeito às diferenças ( linguísticas, culturais, religiosas, étnico-raciais, gênero, entre outros ) tendo a compreensão \textbf{que} essas relações são \textbf{historicamente} construídas. \textbf{Problematização} e desnaturalização das formas de desigualdade, preconceito, discriminação e intolerância no âmbito local, nacional e global. f ) Sociedade, Política e Ética A compreensão da sociedade em \textbf{que} se \textbf{vive}, os problemas e as possibilidades \textbf{que} a cercam são preocupações \textbf{que} todo cidadão deve ter.
Assim, compreender o cenário histórico, socioespacial, cultural, político e econômico subsidia os estudantes a perceberem de forma mais significativa a realidade circundante.
Permite -lhes levar em conta como os indivíduos, no ritmo acelerado de suas experiências cotidianas, adquirem frequentemente \textbf{uma} consciência falsa de suas posições sociais.
Caminhando por meio da consciência dos fatos \textbf{explícitos}, a indiferença social e política se transformam em participação nas questões e debates públicos, estimulando todos no exercício contínuo, dinâmico da cidadania, assegurando -lhes, por meio da participação e da vigilância, seus direitos e deveres inclusos aos princípios constitucionais e de respeito aos direitos humanos.
Essa Unidade Temática permite aos estudantes \textbf{explicitar}, debater, analisar e criticar ideias e dados sobre o fenômeno político, sem esquecer de respeitar os diferentes \textbf{posicionamentos}.
Desse modo, espera -se \textbf{que} os estudantes \textbf{enquanto} protagonistas reconheçam \textbf{que} o debate público - marcado pelo respeito à liberdade, autonomia e consciência crítica — oriente escolhas e fortaleça o exercício da cidadania, o respeito e a responsabilidade a diferentes projetos de vida.
Em se tratando da ética, ressalta -se \textbf{que} no contexto da BNCC, deve ser compreendida como “ juízo de apreciação da conduta humana, necessária para o \textbf{viver} em sociedade, e em cujas bases destacam -se as ideias de justiça, solidariedade e livre-arbítrio [... ] ” ( BRASIL, 2018b, p. 547 ).
O \textbf{que} significa \textbf{dizer} é \textbf{que} ao mencionar a ética, em Ciências Humanas e Sociais Aplicadas, é compreender \textbf{que} há necessidade de considerar o respeito às diferenças entre pessoas e povos, no intuito de combater quaisquer preconceitos \textbf{que} venham a afetar o convívio social.
Nessa \textbf{direção}, os estudantes \textbf{necessitam} experimentar situações reais \textbf{que} oportunizem a eles sentirem na prática o \textbf{que} a \textbf{teoria} \textbf{preconiza}.
Primeiro movimento para a formação de \textbf{sujeitos} protagonistas, condição \textbf{que} desenvolve \textbf{uma} \textbf{percepção} crítica sobre si e o outro, contribuindo para a autonomia tão almejada por todos os profissionais da educação e ainda, com \textbf{compromisso} ético.
A partir do momento em \textbf{que} os estudantes percebem serem capazes de transformar e serem transformados, por meio de suas ações, adquirem competências para atuarem no mundo polarizado e plural.
Portanto, no ensino médio, a área de Ciências Humanas e Sociais Aplicadas tem como desafio desenvolver a capacidade dos estudantes de agir autonomamente, tanto em processos pessoais quanto nos coletivos, preparando -os de modo integral para a compreensão de seu tempo, espaço, relações e articulações com o seu meio.
Nesse sentido, define -se \textbf{que} o papel da escola é proporcionar aos estudantes a formação integral necessária para a convivência com a diversidade, \textbf{que} abarca não somente o espaço educativo, mas todo o âmbito social.
O parágrafo 1º, art. 11, da Resolução CNE/CEB nº 3/2018, orienta \textbf{que} a organização por áreas tem o sentido de tornar mais \textbf{sólidas} as relações entre os saberes e, dessa forma, favorecer a \textbf{contextualização} para a apreensão e a intervenção na realidade.
Na sequência, serão destacados os aspectos específicos constantes dos componentes curriculares dessa área.

% matched lemmas: alicerce, atual, autor, categoria, compromisso, conceitual, contemporâneo, contextualização, detalhar, direção, divisão, dizer, enquanto, equidade, exigir, explicitar, explícito, fronteira, historicamente, hoje, homem, humanos, implicar, imprescindível, influenciar, necessitar, percepção, permear, posicionamento, preconizar, problematização, que, raciocínio, sistematização, sujeito, sólido, teoria, teórico, um, ver, viver
\end{textsample}
